% Options for packages loaded elsewhere
\PassOptionsToPackage{unicode}{hyperref}
\PassOptionsToPackage{hyphens}{url}
\documentclass[
  ignorenonframetext,
]{beamer}
\newif\ifbibliography
\usepackage{pgfpages}
\setbeamertemplate{caption}[numbered]
\setbeamertemplate{caption label separator}{: }
\setbeamercolor{caption name}{fg=normal text.fg}
\beamertemplatenavigationsymbolsempty
% remove section numbering
\setbeamertemplate{part page}{
  \centering
  \begin{beamercolorbox}[sep=16pt,center]{part title}
    \usebeamerfont{part title}\insertpart\par
  \end{beamercolorbox}
}
\setbeamertemplate{section page}{
  \centering
  \begin{beamercolorbox}[sep=12pt,center]{section title}
    \usebeamerfont{section title}\insertsection\par
  \end{beamercolorbox}
}
\setbeamertemplate{subsection page}{
  \centering
  \begin{beamercolorbox}[sep=8pt,center]{subsection title}
    \usebeamerfont{subsection title}\insertsubsection\par
  \end{beamercolorbox}
}
% Prevent slide breaks in the middle of a paragraph
\widowpenalties 1 10000
\raggedbottom
\AtBeginPart{
  \frame{\partpage}
}
\AtBeginSection{
  \ifbibliography
  \else
    \frame{\sectionpage}
  \fi
}
\AtBeginSubsection{
  \frame{\subsectionpage}
}
\usepackage{iftex}
\ifPDFTeX
  \usepackage[T1]{fontenc}
  \usepackage[utf8]{inputenc}
  \usepackage{textcomp} % provide euro and other symbols
\else % if luatex or xetex
  \usepackage{unicode-math} % this also loads fontspec
  \defaultfontfeatures{Scale=MatchLowercase}
  \defaultfontfeatures[\rmfamily]{Ligatures=TeX,Scale=1}
\fi
\usepackage{lmodern}
\ifPDFTeX\else
  % xetex/luatex font selection
\fi
% Use upquote if available, for straight quotes in verbatim environments
\IfFileExists{upquote.sty}{\usepackage{upquote}}{}
\IfFileExists{microtype.sty}{% use microtype if available
  \usepackage[]{microtype}
  \UseMicrotypeSet[protrusion]{basicmath} % disable protrusion for tt fonts
}{}
\makeatletter
\@ifundefined{KOMAClassName}{% if non-KOMA class
  \IfFileExists{parskip.sty}{%
    \usepackage{parskip}
  }{% else
    \setlength{\parindent}{0pt}
    \setlength{\parskip}{6pt plus 2pt minus 1pt}}
}{% if KOMA class
  \KOMAoptions{parskip=half}}
\makeatother
\usepackage{longtable,booktabs,array}
\usepackage{caption}
\captionsetup[table]{skip=6pt}
\newcounter{none} % for unnumbered tables
\usepackage{calc} % for calculating minipage widths
\usepackage{caption}
% Make caption package work with longtable
\makeatletter
\def\fnum@table{\tablename~\thetable}
\makeatother
\setlength{\emergencystretch}{3em} % prevent overfull lines
\providecommand{\tightlist}{%
  \setlength{\itemsep}{0pt}\setlength{\parskip}{0pt}}
% Auto-generated by infrastructure/rendering/_pdf_latex_helpers.py::extract_math_font_preamble.
% Minimal header passed to Pandoc via -H so Beamer slides pick up the same math font
% as the combined-PDF path. Adjust the manuscript's preamble.md to change the font globally.
\usepackage{unicode-math}
\setmathfont{latinmodern-math.otf}

% Slide layout defaults for warning-clean scientific decks.
\usepackage{etoolbox}
\IfFileExists{xurl.sty}{\usepackage{xurl}}{}
\IfFileExists{seqsplit.sty}{\usepackage{seqsplit}}{\newcommand{\seqsplit}[1]{#1}}
\protected\def\breakseq#1{\seqsplit{#1}}
\protected\def\breaktt#1{\begingroup\ttfamily\seqsplit{#1}\endgroup}
\setlength{\emergencystretch}{6em}
\tolerance=5000
\hbadness=10000
\hfuzz=1pt
\setlength{\tabcolsep}{2pt}
\AtBeginEnvironment{longtable}{\tiny\renewcommand{\arraystretch}{0.86}\setlength{\tabcolsep}{1pt}}
\AtBeginEnvironment{tabular}{\tiny\renewcommand{\arraystretch}{0.86}\setlength{\tabcolsep}{1pt}}
\AtBeginEnvironment{equation}{\tiny}
\AtBeginEnvironment{equation*}{\tiny}
\AtBeginEnvironment{align}{\tiny}
\AtBeginEnvironment{align*}{\tiny}
\AtBeginEnvironment{itemize}{\footnotesize}
\AtBeginEnvironment{enumerate}{\footnotesize}
\AtBeginEnvironment{description}{\footnotesize}
\setbeamerfont{caption}{size=\tiny}
\setbeamerfont{caption name}{size=\tiny}
\setbeamerfont{normal text}{size=\small}
\setbeamerfont{frametitle}{size=\small}
\setbeamerfont{section title}{size=\footnotesize}
\setbeamerfont{subsection title}{size=\footnotesize}
\setbeamertemplate{section page}{%
  \centering
  \begin{beamercolorbox}[sep=12pt,center,wd=\paperwidth]{section title}
    \parbox{0.86\paperwidth}{\centering\usebeamerfont{section title}\insertsection\par}
  \end{beamercolorbox}
}
\setbeamertemplate{subsection page}{%
  \centering
  \begin{beamercolorbox}[sep=8pt,center,wd=\paperwidth]{subsection title}
    \parbox{0.86\paperwidth}{\centering\usebeamerfont{subsection title}\insertsubsection\par}
  \end{beamercolorbox}
}
\setlength{\abovecaptionskip}{2pt}
\setlength{\belowcaptionskip}{0pt}

% Natbib and cross-reference fallbacks — slides don't load natbib
% or cleveref, but combined-PDF manuscript prose may emit these
% commands. The fallback renders citations as a bracketed key list
% and cross-references as detokenized labels so slides don't fail on
% undefined control sequences or raw underscores. \providecommand is
% a no-op if packages load later.
\providecommand{\citep}[1]{[#1]}
\providecommand{\citet}[1]{#1}
\providecommand{\citealp}[1]{#1}
\providecommand{\citeauthor}[1]{#1}
\providecommand{\citeyear}[1]{#1}
\providecommand{\cref}[1]{\texttt{\detokenize{#1}}}
\providecommand{\Cref}[1]{\texttt{\detokenize{#1}}}

\newtheorem{proposition}{Proposition}
\newtheorem{hypothesis}{Hypothesis}
\usepackage{bookmark}
\IfFileExists{xurl.sty}{\usepackage{xurl}}{} % add URL line breaks if available
\urlstyle{same}
% fallback for those not using the hyperref driver hyperxmp:
\makeatletter
\@ifundefined{xmpquote}{\newcommand{\xmpquote}[1]{#1}}{}
\makeatother
\hypersetup{
  hidelinks,
  pdfcreator={LaTeX via pandoc}}

\author{\texorpdfstring{}{}}
\date{}

\begin{document}

\section{Methods: Evidence Roles, Curriculum Graph, and Stratum
Design}\label{sec:methodology-scholarship-curriculum-strata}

\subsection{Source-role matrix for citation coverage and source
health}\label{source-role-matrix-for-citation-coverage-and-source-health}

\begin{frame}[fragile,allowframebreaks]{Source-role matrix for citation
coverage and source health}
The manuscript uses a second, non-numeric integrity gate for citations.
The static \breaktt{data/scholarship\_sources.yaml} file maps each cited
BibTeX key to a scholarship lane, source kind, locator, and support
role. \breaktt{src/scholarship/coverage.py} parses the authored
manuscript sections, validates those keys against
\breaktt{manuscript/references.bib}, and rejects any cited key that lacks
a matrix entry or authoritative locator. It also rejects stale
source-matrix rows and uncited BibTeX entries, because those orphans
would otherwise let the coverage artifact imply support that no
manuscript sentence actually uses. The resulting JSON and figure are
generated artifacts, not hand-maintained summaries. The source-health
check adds a second layer over this matrix: DOI and URL locators must be
syntactically authoritative, documentation/software rows must point to
URLs, and each role must explain how the source supports the manuscript
\end{frame}

\begin{frame}[fragile,allowframebreaks]{Source-role matrix for citation
coverage and source health}
rather than merely naming the source.

This gate does not decide whether a paper is true. It prevents weaker
failure modes that ordinary citation checks miss: an undefined BibTeX
key, a citation with no stated role, a new source absent from the
coverage matrix, or a source lane that disappears from the manuscript
without anyone noticing. The scholarship figure is therefore an audit
view over the current manuscript, not a bibliography inventory.

The matrix is also a situation map for the paper's evidence boundary.
Foundational and historical sources explain vocabulary and lineage;
Active-Inference sources define or compare model formulations; critique
and empirical-status sources calibrate generality and application
language; pedagogy and assessment sources motivate ordering and future
study design; platform and software sources establish delivery and
implementation context; and provenance, reproducibility, accessibility,
\end{frame}

\begin{frame}[fragile,allowframebreaks]{Source-role matrix for citation
coverage and source health}
and security sources describe artifact obligations. A source can occupy
more than one lane, but its role is declared at the claim boundary
rather than inferred from prestige, recency, or citation count.
\end{frame}

\subsection{Curriculum graph model, DAG validation, and quality
gates}\label{curriculum-graph-model-dag-validation-and-quality-gates}

\begin{frame}[fragile,allowframebreaks]{Pedagogical quality constraints}
\protect\phantomsection\label{pedagogical-quality-constraints}
\texttt{src/curriculum} defines typed \texttt{Stratum},
\texttt{Subject}, \texttt{Skill}, and \breaktt{AssessmentQuestion}
records, with field names mirroring SkillTree so export is a thin map.
The \texttt{CurriculumGraph} validates structural invariants -- unique
identifiers, no self-loops, no dangling prerequisites, no cycles, and
roots only in the two foundational strata -- and provides deterministic
topological order, per-skill depth levels, and a reachability check from
foundations to applications. Cycle detection uses the standard
three-color depth-first-search pattern from graph algorithms (Cormen et
al. 2022).

\end{frame}

\begin{frame}[fragile,allowframebreaks]{Pedagogical quality constraints}
Structural validity is necessary but not sufficient for learning
quality, so the build also runs \breaktt{src/curriculum/quality.py}. That
gate rejects shallow one-question skills, missing tags, terse subject or
skill descriptions, answer-position bias, and backward prerequisites
that would force later-stratum concepts into earlier learning paths.
Such cross-stratum references are retained as non-blocking related links
rather than exported as SkillTree dependencies. A semantic-prerequisite
gate additionally rejects any skill that is \emph{assessed on} a
load-bearing concept (KL divergence, Markov chains, the Laplace
approximation, the ELBO) without the concept's defining skill in its
transitive prerequisites, flagging pedagogical inversions while
exempting forward references for motivation; a non-empty-scan guard
fails the build if any guarded concept stops appearing in the live
corpus, so the gate cannot silently become vacuous.
\end{frame}

\begin{frame}[fragile,allowframebreaks]{Pedagogical quality constraints}

The newest hardening layer expands continuous dynamics, state-space
probability, rate-distortion information, continuous variational
inference, predictive-coding process theory, hierarchical active
inference, reproducible simulation, and applied evaluation dossiers.
Those additions are grounded in filtering, information bottleneck,
Helmholtz-machine, and generalized-filtering references while remaining
pedagogical rather than production inference engines (Särkkä 2013;
Tishby et al. 1999; Dayan et al. 1995; Friston et al. 2010).

\end{frame}

\begin{frame}[fragile,allowframebreaks]{Pedagogical quality constraints}
The same graph is educationally meaningful only in a narrower sense: it
approximates a knowledge-space style closure condition over authored
skills, where later tasks should not be treated as feasible before their
prerequisites, but it does not infer an individual learner's latent
knowledge state (Doignon and Falmagne 1985; Falmagne and Doignon 2011).
\end{frame}

\subsection{Stratum order as an inspectable historical and pedagogical
design
choice}\label{stratum-order-as-an-inspectable-historical-and-pedagogical-design-choice}

\begin{frame}[allowframebreaks]{Stratum order as an inspectable
historical and pedagogical design choice}
The eight-stratum ordering -- mathematics, probability, information
theory, variational methods, the Free Energy Principle, Active
Inference, computation, and applications -- is an authored design
decision, not a unique decomposition of the domain. An
applications-first path or a thematic split between generative-model
statics and inference/action dynamics could serve different learners.
This release chooses the linear order because the graph can then enforce
a simple local rule: concepts precede the concepts that depend on them
(Bloom 1956; Bloom 1968; Liang et al. 2018). Historical sources motivate
the categories and expose older problems of demonstration, perception,
probability, and ordered instruction, but they do not determine the
graph edges or certify the selected order as historically inevitable
(Aristotle 1901; Ibn al-Haytham 1989; Quintilian 95 AD; Huygens 1657).

\end{frame}

\begin{frame}[allowframebreaks]{Stratum order as an inspectable
historical and pedagogical design choice}
Historical pedagogy sources support this as a bounded history of design
concerns rather than as an efficacy claim: Bacon emphasizes staged
inquiry, Comenius ordered progression, Locke habit and practice, and
Rousseau developmental sequencing (Bacon 1620; Comenius 1657; Locke
1693; Rousseau 1762). Knowledge-space and learning-space theory provide
a formal comparison for prerequisite-constrained mastery states, while
curriculum-aware cognitive diagnosis work shows that a stronger
learner-model claim would require response data and an explicit model of
learner state (Doignon and Falmagne 1985; Falmagne and Doignon 2011; Fu
and Fang 2025). Hugo's Active Inference account of teacher development,
embodied cognition, and education-oriented Active Inference work
motivate keeping the order revisable and uncertainty-sensitive, but this
release does not evaluate teacher development, infer learner states, or
\end{frame}

\begin{frame}[allowframebreaks]{Stratum order as an inspectable
historical and pedagogical design choice}
implement an adaptive diagnostic engine (Hugo 2026; Clark 2016; Di Paolo
et al. 2024).

Microlearning and gamification sources motivate modular delivery while
warning that implementation context, satisfaction, and measured
effectiveness can diverge (Septiani and Rosmansyah 2021; Monib et al.
2025; Rof et al. 2024). Active-learning, epistemic-curiosity,
cognitive-engagement, retrieval-practice, and
evidence-centered-assessment sources therefore inform the design of
interaction and quiz tasks; they do not substitute for data from this
curriculum or establish durable competence (Freeman et al. 2014; Emin et
al. 2026; Chi and Wylie 2014; Roediger and Karpicke 2006; Mislevy et al.
2003). Constructive alignment, mastery learning, formative assessment,
and feedback provide the local rule that objectives, activities, and
quiz gates should be co-specified and used as revision checkpoints
(Biggs 1996; Bloom 1968; Black and Wiliam 1998; Shute 2008). The
\end{frame}

\begin{frame}[allowframebreaks]{Stratum order as an inspectable
historical and pedagogical design choice}
resulting ordering is a parameter to inspect and revise, not a
theoretical or empirically tested claim about how the field must be
structured.
\end{frame}

\end{document}
